% A consequence of the modular-representation compression argument in
% Schopieray, arXiv:2609.35287v1, Lemma 4.1 and proof of Theorem 5.1.
% This file deliberately does not use the paper's identification of
% [K_1:Q(D)] with a dimensional grading for a non-pseudounitary category.

\begin{proposition}[Compression over all constituents of $I(\mathbf1)$]
Let $\mathcal C$ be a spherical fusion category, $D=\dim(\mathcal C)$,
$K=\mathbb Q(\dim X:X\in\operatorname{Irr}(\mathcal C))$, and
$G=\operatorname{Gal}(K/\mathbb Q)$.  For an irreducible representation
$E$ of $K_0(\mathcal C)\otimes\mathbb C$, write $f_E$ for its formal
codegree.  Then $f_E\in K$, and for every character $\chi$ of $G$,
\begin{equation}\label{eq:full-I1-compression}
 \sum_E\left|\sum_{\sigma\in G}
       \chi(\sigma)\sigma(f_E^{-1})\right|^2
       \leq \frac{|G|}{\mathfrak f(\chi)}.
\end{equation}
There is no factor $\dim E$ in the outer sum.
\end{proposition}

\begin{proof}
Let $A_E$ be the simple constituent of $I(\mathbf1)$ attached to $E$.
The induction theorem gives
$\dim A_E=D/f_E$, $\theta_{A_E}=1$, and distinct $E$ give distinct
$A_E$.  Since the forgetful image of every object of the center has
dimension in $K$, this also proves $f_E\in K$.
Write $s=S/D$ for the normalized center matrix and $P$ for orthogonal
projection onto its twist-one subspace.  For the dimension character and
its Galois conjugates let $A_\sigma$ denote the corresponding constituents
of $I(\mathbf1)$, and set
$w_\chi=\sum_{\sigma\in G}\chi(\sigma)e_{A_\sigma}$.
The Galois compatibility of the modular representation gives
$s_{A_\sigma,Y}=\sigma(\dim Y/D)$, and $w_\chi$ is fixed by the
unipotent subgroup and has torus character $\chi$.  The local compression
estimate for $SL_2(\mathbb Z/N)$, combined over prime powers, gives
$\|Ps w_\chi\|^2\leq\|w_\chi\|^2/\mathfrak f(\chi)
=|G|/\mathfrak f(\chi)$.
Its $A_E$ coordinate is exactly the inner sum in
\eqref{eq:full-I1-compression}.  Retaining these coordinates proves the
claim.  The multiplicity $[I(\mathbf1):A_E]=\dim E$ does not repeat a
coordinate of the modular representation.
\end{proof}

\begin{corollary}[Quadratic Casimir obstruction]\label{cor:quadratic-Casimir}
Let $R$ be a fusion ring, let $N_i$ be its regular fusion matrices, and let
$B=\sum_iN_iN_i^{\mathsf T}$.  Fix a real quadratic field $L$ of
fundamental discriminant $q$.  List the irreducible quadratic factors of
$\det(xI-B)$ whose roots generate $L$ as
\[
 \prod_j(x^2-t_jx+n_j)^{m_j},
 \qquad
 Q_L(R):=q\sum_jm_j\frac{t_j^2-4n_j}{n_j^2}.
\]
Assume that at least one such factor occurs.  If $R$ has a spherical categorification, then
$Q_L(R)\leq1$.  More precisely, if $K$ is its dimension field, then
$L\subseteq K$ and $[K:L]Q_L(R)\leq1$.
\end{corollary}

\begin{proof}
The Casimir eigenvalue on a simple representation $E$ is
$\lambda_E=(\dim E)f_E$, with regular multiplicity $(\dim E)^2$.
In particular $L\subseteq K$.  Put $e'=[K:L]$ and pull back to
$\operatorname{Gal}(K/\mathbb Q)$ the nontrivial quadratic character of
$L$, whose conductor is $q$.  If $\lambda_E\in L$, its contribution to
\eqref{eq:full-I1-compression} is
$e'^2(\dim E)^2(\lambda_E^{-1}-\overline{\lambda_E}^{-1})^2$.
Regular multiplicities therefore account for the squares of representation
degrees, including when $R$ is noncommutative.  Each displayed factor
contributes $2e'^2m_j(t_j^2-4n_j)/n_j^2$.  The upper bound is $2e'/q$,
which gives the asserted inequality after cancellation.
\end{proof}

\begin{corollary}[Boundary constraint for an arbitrary categorification]
If $R$ has any fusion-category categorification, then $Q_L(R)\leq4$.
More generally, if $F$ is a field generated by Casimir eigenvalues and
$L\subseteq F$, then $[F:L]Q_L(R)\leq4$.
If $Q_L(R)=4$, its sphericalization has dimension field exactly $L$,
and every formal codegree of its $\varepsilon=-1$ sector is rational.
In particular its actual spherical global dimension is rational.
\end{corollary}

\begin{proof}
The sphericalization has two lifts of every simple object and a central
invertible $\varepsilon$ of order two and dimension $-1$.
The $\varepsilon=1$ quotient of its fusion algebra is the original fusion
algebra; in this sector its Casimir is $2B$.  Thus the factors supplied by
this sector contribute $Q_L(R)/4$ to the preceding proof.  Its dimension
field contains $L$, so $[K:L]Q_L(R)/4\leq1$.
Every Casimir eigenvalue is a rational integer multiple of a formal
codegree, so also $F\subseteq K$, which proves the stronger bound.
At equality $Q_L(R)=4$, necessarily $K=L$, and the known sector already
exhausts the squared-norm bound.  All other $I(\mathbf1)$ coordinates must
vanish under the nontrivial Fourier character of $L$, so their formal
codegrees are rational.  The spherical dimension character belongs to the
$\varepsilon=-1$ sector, and its formal codegree is the global dimension.
\end{proof}

\paragraph{Scope caution.}
The last corollary does not imply a nontrivial grading of the original
ring.  For arbitrary spherical categories, rational global dimension
does not force rational squared dimensions of simples: the product of a
Fibonacci category and its Yang--Lee conjugate has global dimension $5$,
dimension field $\mathbb Q(\sqrt5)$, and trivial universal grading.
Thus the Frobenius--Perron dimensional-grading theorem must not be applied
to an arbitrary spherical dimension character.

\paragraph{Two rank-eight applications.}
Both $R8\_2\_318$ and $R8\_2\_735$ have Casimir characteristic polynomial
\[
 (x-16)^3(x-8)(x^2-192x+1024)(x^2-16x+32).
\]
The two quadratic factors lie in $\mathbb Q(\sqrt2)$, of conductor $8$,
and give
\[
 Q_{\mathbb Q(\sqrt2)}
 =8\left(\frac{192^2-4\cdot1024}{1024^2}
          +\frac{16^2-4\cdot32}{32^2}\right)=\frac54>1.
\]
Consequently neither ring has a spherical categorification.  The explicit
eight characters show that the only characters nonzero on every basis
element are the Frobenius--Perron character and its real quadratic
conjugate.  Both are invariant under duality.  Since simple pivotal
dimensions are nonzero, every pivotal categorification would therefore
be spherical.  Neither ring has a pivotal categorification.
This leaves arbitrary nonpivotal categorifications unresolved.
